
% \begin{table}[H]
%     \caption{Reclassification scheme for aligning NLCD land cover classes with our desired classification scheme.}\label{tab:nlcd_reclass}
%     \vspace{0.25cm}
%     \centering
%     \resizebox{\textwidth}{!}{
%         \begin{tabular}{@{}llll@{}}
%             \toprule
%             NLCD Class Index & NLCD Class                  & Mississippi Class Index & Mississippi Class                      \\ \midrule
%             11               & Open Water                  & 1                       & Open Water                             \\ \smallrule
%             \makecell[l]{23 \\ 24} & \makecell[l]{Developed, Medium Intensity \\ Developed, High Intensity} & 2                       & Impervious Structures                   \\ \smallrule
%             \makecell[l]{21 \\ 22} & \makecell[l]{Developed, Open Space \\ Developed, Low Intensity} & 3                       & Impervious Surfaces                     \\ \smallrule
%             31               & Barren Land                 & 4                       & Barren                                  \\ \smallrule
%             \makecell[l]{41 \\ 42 \\ 43 \\ 90} & \makecell[l]{Deciduous Forest \\ Evergreen Forest \\ Mixed Forest \\ Woody Wetlands} & 5                       & Forest                                  \\ \smallrule
%             \makecell[l]{52 \\ 71 \\ 81 \\ 95} & \makecell[l]{Shrub/Scrub \\ Grassland/Herbaceous \\ Pasture/Hay \\ Emergent Herbaceous Wetlands} & 6                       & Herbaceous                              \\ \smallrule
%             82               & Cultivated Crops            & 7                       & Cultivated Crops                       \\ \smallrule
%             250              & NoData                      & 0                       & No Data                                \\ \bottomrule
%         \end{tabular}
%     }
% \end{table}

\begin{table}[H]
    \caption{Reclassification scheme for aligning NLCD land cover classes with our desired classification scheme.}\label{tab:nlcd_reclass}
    \centering
    \begin{tblr}{
        colspec={Q[c,m] X[l,m] Q[c,m] Q[l,m]},
        row{1} = {font=\bfseries},
        hlines = {0.25pt, dashed},
        hline{1,2,10} = {0pt, solid},
        % booktabs
    }
        \toprule
        {NLCD Class\\Index} & NLCD Class & {Target Class\\Index} & Target Class \\
        \midrule
        11 & Open Water & 1 & Open Water \\
        {23 \\ 24} & {Developed, Medium Intensity \\ Developed, High Intensity} & 2 & Impervious Structures \\
        {21 \\ 22} & {Developed, Open Space \\ Developed, Low Intensity} & 3 & Impervious Surfaces \\
        31 & Barren Land & 4 & Barren \\
        {41 \\ 42 \\ 43 \\ 90} & {Deciduous Forest \\ Evergreen Forest \\ Mixed Forest \\ Woody Wetlands} & 5 & Forest \\
        {52 \\ 71 \\ 81 \\ 95} & {Shrub/Scrub \\ Grassland/Herbaceous \\ Pasture/Hay \\ Emergent Herbaceous Wetlands} & 6 & Herbaceous \\
        82 & Cultivated Crops & 7 & Cultivated Crops \\
        250 & No Data & 0 & No Data \\
        \bottomrule
    \end{tblr}
\end{table}